\section{Pendahuluan}
\label{sec:pendahuluan}

\subsection{Latar Belakang}

Penanggulangan bencana melibatkan lembaga-lembaga yang bekerja secara otonom,
memiliki sumber data, format informasi, dan karakteristik infrastruktur yang
berbeda. Dalam skenario tugas ini, BMKG menyediakan informasi kejadian seismik
dan potensi tsunami, sedangkan PVMBG menyediakan laporan aktivitas vulkanik.
BNPB membutuhkan kedua aliran informasi tersebut untuk membangun gambaran
situasi yang terpadu dan menyajikannya kepada pihak yang berwenang. Namun,
perbedaan format data, latensi layanan, dan batas kepercayaan antarlembaga
membuat integrasi langsung menjadi tidak memadai.

Masalah tersebut tidak hanya berkaitan dengan pertukaran data. Ketika satu
layanan sumber melambat atau gagal, layanan lain tetap dapat menghasilkan data
yang bernilai bagi pengguna. Demikian pula, lonjakan permintaan pada saat status
bencana meningkat tidak boleh membuat seluruh sistem berhenti merespons.
Sistem juga perlu membedakan hak akses antara Media, Tim Lapangan, dan BNPB
Internal agar data mentah seperti koordinat presisi serta atribut spesifik sumber
tidak bocor kepada pihak yang hanya membutuhkan ringkasan.

Oleh sebab itu, proyek ini membangun proof-of-concept platform koordinasi
kebencanaan BNPB sebagai aplikasi terdistribusi. Platform menyatukan data
heterogen dari BMKG dan PVMBG melalui komunikasi jaringan, mengubahnya ke
format kanonik \textit{HazardEvent}, menyimpannya pada \textit{Canonical Store},
dan mendistribusikan informasi tersebut melalui API serta mekanisme
\textit{publish-subscribe}.

\subsection{Tujuan}

Tujuan utama pengembangan sistem adalah membuktikan bahwa beberapa layanan
otonom dapat diintegrasikan menjadi satu platform koordinasi tanpa menghilangkan
batas kepemilikan dan karakteristik masing-masing layanan. Secara khusus,
Milestone 1 bertujuan untuk:

\begin{enumerate}
  \item menyatukan data seismik dan vulkanik dari dua instansi ke dalam format
  kanonik yang konsisten serta tetap menerima perubahan skema aditif;
  \item menjaga ketersediaan data penting melalui pemrosesan asinkron,
  pembatasan konkurensi, dan degradasi layanan yang eksplisit ketika salah satu
  sumber tidak tersedia;
  \item menerapkan pemisahan domain kepercayaan untuk komunikasi antarlayanan
  dan akses client downstream;
  \item memisahkan komponen ke dalam layanan dan container yang dapat dijalankan
  secara independen, dengan penyimpanan yang mendukung atribut dinamis; serta
  \item mendistribusikan event baru ke beberapa consumer melalui message broker
  tanpa membuat Aggregator bergantung langsung pada consumer tersebut.
\end{enumerate}

\subsection{Ruang Lingkup Milestone 1}

Ruang lingkup implementasi dibatasi pada dua instansi sumber dan dua jenis
bencana. BMKG disimulasikan sebagai sumber \textit{SeismicEvent} dan
\textit{TsunamiWarning}; PVMBG disimulasikan sebagai sumber
\textit{VolcanicReport}. Kedua sumber bersifat pasif dan hanya merespons polling
berkala dari Aggregator BNPB. Data dari sumber tidak dibagikan melalui pemanggilan
fungsi langsung atau library \textit{business logic} bersama, melainkan melalui
API jaringan sesuai kontrak layanan.

Komponen BNPB dalam Milestone 1 meliputi Aggregator atau gateway, Canonical
Store, Client-Facing API, Auth/Token Service, message broker, dan sedikitnya dua
event consumer. Aggregator melakukan polling, pemetaan, penyimpanan, serta
publikasi event baru. Canonical Store hanya dimiliki dan diakses langsung oleh
Aggregator; layanan lain memperoleh data melalui API Aggregator. Client-Facing
API bersifat \textit{read-only} dan menerapkan hak akses berdasarkan scope token.

Hak akses client dibagi menjadi tiga tingkat. Media hanya menerima field
berklasifikasi Ringkasan. Tim Lapangan memperoleh data lengkap untuk kebutuhan
operasional, sedangkan BNPB Internal memiliki hak akses tertinggi. Access token
bersifat berumur pendek dan harus dapat diperbarui melalui alur \textit{refresh}
tanpa login ulang. Batasan ini dibuat agar keamanan, ketahanan, dan distribusi
data dapat diuji secara nyata pada arsitektur yang sama.

\subsection{Batasan Sistem}

Milestone 1 belum mencakup seluruh kebutuhan sistem kebencanaan nasional.
Sistem tidak menangani jalur pelaporan dari lapangan ke BNPB, perubahan data oleh
client downstream, maupun sumber bencana lain seperti kebakaran hutan dan
hidrometeorologi. Fokusnya adalah fondasi interoperabilitas, availability,
keamanan, fleksibilitas storage, dan distribusi event pada skala dasar.

Pengembangan untuk Milestone 2 berada di luar cakupan laporan ini. Tahap
berikutnya akan menguji fondasi yang sama dalam kondisi yang lebih skalabel,
termasuk banyak instance layanan, volume data yang meningkat drastis, dan
kegagalan beberapa komponen secara bersamaan.

\subsection{Struktur Laporan}

Bab berikutnya menjelaskan arsitektur sistem, batas tiap service, serta teknologi
yang digunakan. Setelah itu, setiap problem resmi Milestone 1 dibahas secara
terpisah, mencakup rancangan solusi, alternatif dan trade-off, batasan yang
disadari, serta bukti penyelesaian. Lampiran digunakan untuk menyimpan bukti
pengujian tambahan seperti tangkapan layar, log ber-correlation ID, dan hasil
pengukuran.
